% SOURCE_STATEMENT: "Is the so-called four-qubit $W$ state in the LOCC dense-codeable class?"
% Preparation and final self-audit: CONTRIBUTING.md prompts C and D; source/audit artifacts retained in the local physics_statements_2026_10_11 batch.
% !TeX program = lualatex
\documentclass[11pt]{article}
\usepackage[margin=25mm]{geometry}
\usepackage{fontspec}
\setmainfont{Latin Modern Roman}
\usepackage{amsmath,amssymb,mathtools}
\usepackage{unicode-math}
\setmathfont{Latin Modern Math}
\usepackage{xurl,hyperref}
\hypersetup{colorlinks=true,urlcolor=blue}
\setcounter{secnumdepth}{0}
\setlength{\parindent}{0pt}
\setlength{\parskip}{5pt}
\setlength{\emergencystretch}{3em}
\newcommand{\sref}[2]{\hyperlink{source:#1:#2}{[S#2]}}
\newcommand{\cataloguescope}[1]{\paragraph{Recorded target.}#1}
\newcommand{\R}{\mathbb R}
\newcommand{\C}{\mathbb C}
\newcommand{\N}{\mathbb N}
\newcommand{\Z}{\mathbb Z}
\newcommand{\Q}{\mathbb Q}
\newcommand{\F}{\mathbb F}
\newcommand{\eps}{\varepsilon}
\begin{document}
\section{LOCC Dense-Codeability of the Four-Qubit W State}\label{statement}

\subsection{Definitions and mathematical statement}

\paragraph{Shared notation.}$\mathbb N=\{1,2,\ldots\}$; the standard number fields have their usual meanings. All additional conventions are stated locally.


\paragraph{Definitions and conventions.} A qubit has Hilbert space $\mathbb C^2$, and the normalized four-qubit state is $|W_4\rangle=\tfrac12(|1000\rangle+|0100\rangle+|0010\rangle+|0001\rangle)$. In the report's two-receiver protocol, two sending qubits are locally encoded and transmitted to their assigned receivers. The two receivers then decode by LOCC. A state is LOCC dense-codeable when the achievable classical communication rate exceeds the rate available from the transmitted systems without the shared entangled state; with two sending qubits the benchmark is two bits per use.

\paragraph{Source scope and qualifications.} This is the report's two-receiver LOCC question, not global decoding, a single-receiver protocol, or a teleportation question. Receiver partition and the report's encoding class remain part of the target.

\paragraph{Catalogue mathematical target.}
\paragraph{Catalogue wording retained for traceability.}
Is the so-called four-qubit $W$ state in the LOCC dense-codeable class?

\paragraph{Saved literature context.} The supplied catalogue's assessment dated 2026-08-21 records: The cited dense-coding literature develops capacity bounds and examples but does not settle whether the standard four-qubit W state is LOCC dense codeable under the report's two-receiver protocol. Later W-state communication results address related protocols, not this exact question. The assessment is retained as dated provenance; it is not a new comprehensive literature or open-status certification.
\subsection{Short English statement}
Can the four-qubit W state provide a dense-coding advantage when its two receivers communicate only classically?
\subsection{Sources}
\begin{itemize}
\item[\hypertarget{source:statement:1}{S1}] ULAM.AI and the UnsolvedMath Contributors, UnsolvedMath, catalogue entry. Catalogue curation under CC BY 4.0; definitions and English exposition added locally.
\url{https://huggingface.co/datasets/ulamai/UnsolvedMath}.
\item[\hypertarget{source:statement:2}{S2}] Entanglement and Decoherence: Mathematics and Physics of Quantum Information and Computation (2005), Oberwolfach Reports, DOI 10.4171/owr/2005/04, pp. 185--255; official report source and mathematical context.
\url{https://doi.org/10.4171/owr/2005/04}.
\item[\hypertarget{source:statement:3}{S3}] Official EMS Press PDF of the same Oberwolfach report; PDF pages 19--21, printed pages 203--205.
\url{https://ems.press/content/serial-article-files/45978}.
\item[\hypertarget{source:statement:4}{S4}] Bruss et al., Dense coding with multipartite quantum states; reference in the saved dated catalogue literature review.
\url{https://arxiv.org/abs/quant-ph/0507146}.
\item[\hypertarget{source:statement:5}{S5}] Yeo--Chua, teleportation and dense coding with genuine multipartite entanglement; reference in the saved dated catalogue literature review.
\url{https://doi.org/10.1103/PhysRevLett.96.060502}.
\end{itemize}
\label{end:statement}
\end{document}
